\item \textbf{Problema de repaso.}
Un bloque de $0.250\text{ kg}$ que descansa sobre una superficie horizontal sin fricción se une a un resorte cuya constante de fuerza es $83.8\text{ N/m}$ como en la figura P15.15. Una fuerza horizontal $\vec{F}$ hace que el resorte se estire una distancia de $5.46\text{ cm}$ desde su posición de equilibrio.
\begin{enumerate}
    \item Encuentre la magnitud de $\vec{F}$.
    \item ¿Cuál es la energía total almacenada en el sistema cuando se estira el resorte?
    \item Encuentre la magnitud de la aceleración del bloque justo después de que se retira la fuerza aplicada.
    \item Encuentre la rapidez del bloque cuando éste primero alcanza la posición de equilibrio.
    \item Si la superficie no es sin fricción, pero el bloque todavía alcanza la posición de equilibrio, ¿sería su respuesta a la parte (4) mayor o menor?
    \item ¿Qué otra información necesitaría saber para encontrar la respuesta real del inciso (4) en este caso?
    \item ¿Cuál es el valor más grande del coeficiente de fricción que permita al bloque alcanzar la posición de equilibrio?
\end{enumerate}